% §14.3 Space Analysis
%
% Cross-reference resolution status:
%
% DEFINED here:
%   \label{sec:space-analysis}                   -- this section
%   \label{thm:bhf-optimality}                   -- the BHF optimality theorem
%
% ALREADY DEFINED (chapter 5):
%   \Cref{def:bloom-filter}                      -- chapter 5 §5.2
%
% ALREADY DEFINED (chapter 8):
%   \Cref{thm:entropy-floor}                     -- chapter 8 §8.3
%   \Cref{prop:bloom-overhead}                   -- chapter 8 §8.4
%
% ALREADY DEFINED (chapter 14, earlier sections):
%   \Cref{sec:optimality-question}               -- §14.1
%   \Cref{def:bhf}                               -- §14.2
%   \Cref{sec:bhf-construction}                  -- §14.2
%   \eqref{eq:bhf-r}                             -- §14.2 fingerprint width
%
% FORWARD REFS within chapter 14:
%   \Cref{sec:fpr-analysis}                      -- §14.4 (Task 5)
%   \Cref{sec:construction-time}                 -- §14.5 (Task 6)
%
% §14.3 computes the BHF's bit budget: n*r for the fingerprints,
% sublinear for the perfect hash; states the optimality theorem
% (BHF matches the entropy floor up to integer rounding plus o(1));
% compares against floor and Bloom in a table at n=10^6 with three
% fpr targets; closes with the ceil(.) rounding discussion.

\section{Space Analysis}
\label{sec:space-analysis}

The BHF stores two objects: the perfect hash $\phi$ and the
fingerprint table $F$. The total bit budget is the sum of the two,
and the section's task is to compare that sum to the entropy floor
of \Cref{thm:entropy-floor}. The fingerprint table accounts for the
dominant term and matches the floor up to the integer rounding of
the fingerprint width $r$. The perfect-hash representation accounts
for a sublinear correction that vanishes asymptotically. Putting the
two together gives the chapter's headline result: the BHF achieves
the floor up to one bit of rounding plus a vanishing additive term.
The Bloom filter's $\log_2 e$ multiplicative overhead is gone.

\subsection*{Fingerprint storage}

The fingerprint table $F : [n] \to \{0, 1\}^r$ holds $n$ slots, each
$r$ bits wide. The total bit cost is
\begin{equation}
\label{eq:bhf-fingerprint-bits}
  L_{F} \;=\; n \cdot r \;=\; n \cdot \lceil -\log_2 \fprate \rceil
\end{equation}
bits, where $r = \lceil -\log_2 \fprate \rceil$ is the width from
\eqref{eq:bhf-r}. The dependence on the target rate is logarithmic:
each extra order of magnitude in $1/\fprate$ adds about $3.32$ bits
to $r$, which is to say one extra base-$10$ digit of precision costs
about $3.32$ bits per stored element. The dependence on the set
size is linear in $n$. The fingerprint term scales with both
quantities the designer controls, and scales in the way the floor
predicts.

The integer rounding in $r$ contributes between $0$ and $1$ extra
bits per element compared to the real-valued floor $-\log_2 \fprate$.
The rounding cost is bounded above by one bit per element regardless
of $\fprate$: for any real $x > 0$, $\lceil x \rceil - x \in [0, 1)$.
When $-\log_2 \fprate$ is close to an integer, the rounding is small;
when it is at the midpoint between two integers, the rounding is
about $0.5$ bits; in the worst case it approaches $1$ bit. The
average rounding over a uniformly distributed $\fprate$ on a
logarithmic scale is $0.5$ bits per element.

For comparison, the standard Bloom filter at the same $\fprate$ uses
$\log_2(1/\fprate) \cdot \log_2 e \approx 1.4427 \cdot \log_2(1/\fprate)$
bits per element (\Cref{prop:bloom-overhead}). The Bloom overhead is
multiplicative and grows linearly with $\log_2(1/\fprate)$. The BHF
overhead is additive and bounded by one bit. As $\fprate \to 0$, the
BHF budget approaches the floor; the Bloom filter budget remains
$44\%$ above the floor for every $\fprate$. The contrast is the
chapter's headline.

To see why the fingerprint term takes this clean form, note that
\eqref{eq:bhf-fingerprint-bits} is the only nontrivial space the
fingerprint table consumes. The table has exactly $n$ slots because
the perfect hash bijects $A$ onto $[n]$; each slot stores exactly
$r$ bits because the fingerprint hash $g$ maps into $\{0, 1\}^r$.
There is no padding, no overhead per slot beyond the fingerprint
itself, and no addressing structure inside the table other than the
indexing by $\phi$. The table is a flat array, addressed
deterministically, and the bit cost is the sum of the per-slot
widths. This is the densest possible representation of a uniform
fingerprint per element. The Bloom filter, by contrast, must spread
each insertion across $k$ bits to handle collisions, and the bit
cost includes both the per-insertion spread and the per-element
load factor that the optimal-$k$ analysis selects. The BHF avoids
both penalties by handling collisions at build time and storing
exactly the fingerprint each element needs.

\subsection*{Perfect-hash storage}

The second term in the BHF's bit budget is the space the perfect
hash $\phi$ occupies. A perfect hash function from $A$ to $[n]$ is
itself a data structure that must be stored alongside the
fingerprint table; queries need access to $\phi$ to address the
correct slot. The space cost depends on which perfect-hash
construction one uses.

Modern perfect-hash constructions achieve very low per-element
overhead. The BDZ family of Belazzougui, Botelho, and Dietzfelbinger
(2009; the bibliographic entry appears in \S14.7) achieves
\begin{equation}
\label{eq:bdz-bound}
  L_{\phi} \;=\; O\bigl( n \log \log \log n / \log n \bigr)
\end{equation}
bits in total, which is $o(n)$. The per-element overhead $L_{\phi} /
n = O(\log \log \log n / \log n)$ tends to zero as $n$ grows.
Concrete constants matter at finite $n$, but the asymptotic
behaviour is clear: the perfect-hash term is sublinear in $n$, hence
sublinear in $n r$, hence asymptotically negligible compared to the
fingerprint term \eqref{eq:bhf-fingerprint-bits}.

Earlier constructions (Fredman-Komlos-Szemeredi, Hagerup-Tholey, and
the displacement-based families) all achieve at most $O(n)$ bits
total, with leading constants in the range $2$ to $4$ bits per
element. At realistic scales ($n$ between $10^3$ and $10^9$), the
BDZ-style construction occupies on the order of $2$ to $3$ bits per
element. For the present accounting we take the asymptotic
\eqref{eq:bdz-bound} as the relevant bound and treat the constant
as a small additive correction that the BHF inherits from whichever
perfect-hash family is in use.

The key fact is that the perfect-hash overhead per element vanishes
as $n$ grows. The BHF is therefore an additive-overhead
construction: total bits per element are $r + o(1)$ as $n \to
\infty$. The Bloom filter is a multiplicative-overhead construction:
total bits per element are $r' \cdot \log_2 e$ for any choice of
$r' = \log_2(1/\fprate)$. The two regimes differ in kind, not
merely in constant. As the application scales up, the BHF approaches
the floor; the Bloom filter sits a fixed multiplicative distance
above it.

A quick concrete check: at $n = 10^6$, the asymptotic bound
\eqref{eq:bdz-bound} evaluates to $L_\phi / n \approx O(\log \log \log
10^6 / \log 10^6) = O(\log \log 6 / 20)$, which the asymptotic
constants render at roughly $0.2$ to $0.5$ bits per element in
practical BDZ implementations. The dominant fingerprint term is $r$
bits per element, ranging from $7$ at $\fprate = 10^{-2}$ to $20$ at
$\fprate = 10^{-6}$. The perfect-hash overhead is therefore at most
a few percent of the fingerprint cost at $n = 10^6$, and the
percentage shrinks further at larger $n$. For $n = 10^9$ the
overhead is below $0.1$ bits per element. At the scales BHFs are
deployed, the perfect-hash term is a rounding correction to the
fingerprint term, not a separate concern.

\subsection*{Total bits per element}

Combining the two terms gives the BHF's total bit cost:
\begin{equation}
\label{eq:bhf-total-bits}
  L_{\text{BHF}} \;=\; L_F + L_\phi
  \;=\; n \cdot \lceil -\log_2 \fprate \rceil
  \;+\; O\bigl( n \log \log \log n / \log n \bigr)
\end{equation}
bits. Dividing by $n$,
\begin{equation}
\label{eq:bhf-per-element}
  L_{\text{BHF}} / n
  \;=\; \lceil -\log_2 \fprate \rceil
  \;+\; O(\log \log \log n / \log n).
\end{equation}
The first term is bounded between $-\log_2 \fprate$ and
$-\log_2 \fprate + 1$ by integer rounding. The second term is
$o(1)$. As $n \to \infty$ with $\fprate$ fixed, the BHF's per-element
bit cost approaches the integer rounding of the floor.

\begin{theorem}[BHF optimality]
\label{thm:bhf-optimality}
The Bernoulli Hash Function of \Cref{def:bhf} on a designated set
$A$ with $|A| = n$ and target false-positive rate $\fprate \in (0,
1)$ uses
\[
  L_{\text{BHF}} / n
  \;=\; \lceil -\log_2 \fprate \rceil
  \;+\; o(1)
  \;=\; -\log_2 \fprate \;+\; O(1)
\]
bits per element. The construction achieves the entropy floor of
\Cref{thm:entropy-floor} up to one bit of integer rounding plus a
sublinear correction. The multiplicative overhead $\log_2 e
\approx 1.4427$ that the standard Bloom filter pays
(\Cref{prop:bloom-overhead}) is absent: the BHF's overhead is
additive and bounded by a constant, not multiplicative and growing
with $-\log_2 \fprate$.
\end{theorem}

The theorem is the constructive complement to
\Cref{thm:entropy-floor}. The floor theorem said that no FNR-zero
representation can use fewer than $-\log_2 \fprate$ bits per element
asymptotically; the present theorem says that a specific
representation, the BHF, uses no more than that plus an additive
constant. The combination characterizes the FNR-zero family
tightly: the per-element cost lies in $[-\log_2 \fprate + o(1),
-\log_2 \fprate + 1 + o(1)]$ for every sufficient construction, and
the BHF lives at the top of that interval. The remaining slack is
the unavoidable integer-rounding cost of representing the
fingerprint as a whole number of bits.

\subsection*{Comparison table}

Fix $n = 10^6$ items. \Cref{tab:bhf-comparison} compares three
constructions at three target rates: the entropy floor
$-\log_2 \fprate$ from \Cref{thm:entropy-floor}, the standard Bloom
filter at $\log_2(1/\fprate) \cdot \log_2 e$ bits per element from
\Cref{prop:bloom-overhead}, and the BHF at $\lceil -\log_2 \fprate
\rceil$ bits per element from \Cref{thm:bhf-optimality} (ignoring
the sublinear perfect-hash correction, which contributes a fraction
of a bit at $n = 10^6$).

\begin{table}[h]
\centering
\caption{Bits per element at $n = 10^6$. Floor is the lower bound of
\Cref{thm:entropy-floor}; Bloom is the standard construction at
\Cref{prop:bloom-overhead}; BHF is the optimum of
\Cref{thm:bhf-optimality}, ignoring sublinear perfect-hash overhead.}
\label{tab:bhf-comparison}
\begin{tabular}{l|rrr}
\hline
  $\fprate$  & Floor & Bloom & BHF \\
\hline
  $10^{-2}$  & $6.64$  & $9.59$  & $7$  \\
  $10^{-3}$  & $9.97$  & $14.38$ & $10$ \\
  $10^{-6}$  & $19.93$ & $28.76$ & $20$ \\
\hline
\end{tabular}
\end{table}

The Bloom filter pays a multiplicative overhead of $\log_2 e
\approx 1.44$ at every row; the BHF pays an additive overhead of
$\lceil -\log_2 \fprate \rceil - (-\log_2 \fprate)$, which is at
most one bit and is concretely $0.36, 0.03, 0.07$ bits at the three
rows. The absolute Bloom-versus-BHF gap grows with $-\log_2 \fprate$
(from $2.59$ bits per element at $\fprate = 10^{-2}$ to $8.76$ bits
per element at $\fprate = 10^{-6}$); the absolute BHF-versus-floor
gap stays below one bit per element regardless of the target rate.
At $n = 10^6$ and $\fprate = 10^{-6}$, the absolute savings of BHF
over Bloom is $8.76 \cdot 10^6$ bits, or about $1.1$ megabytes.

\subsection*{The \texorpdfstring{$\lceil \cdot \rceil$}{ceiling} rounding}

The BHF's per-element budget is $\lceil -\log_2 \fprate \rceil$
rather than $-\log_2 \fprate$ because the fingerprint width $r$ is
an integer: there is no way to store a fractional bit in a
fingerprint table addressed by the perfect hash. The integer
rounding contributes between $0$ and $1$ bits per element, and the
contribution is wasted in the sense that the FPR achieved at integer
$r$ is $2^{-r}$, which is at most the target $\fprate$ but is
generally strictly smaller. The BHF over-achieves the FPR target by
the rounding-induced slack.

To quantify the slack, fix a target $\fprate$ and write $-\log_2
\fprate = k + \delta$ where $k$ is an integer and $\delta \in [0, 1)$
is the fractional part. The fingerprint width is $r = \lceil k +
\delta \rceil$, which equals $k$ when $\delta = 0$ and equals $k + 1$
when $\delta > 0$. The achieved FPR is $2^{-r} = 2^{-k-1} =
\fprate / 2^{1 - \delta}$ for $\delta > 0$, which is strictly
smaller than $\fprate$ by a factor between $1$ (at $\delta \to 1$)
and $2$ (at $\delta \to 0^+$). The BHF's FPR is therefore on average
about half the target rate, the price paid for integer rounding.

The wasted bits can in principle be recovered by storing the
fingerprint at fractional precision: arithmetic coding or block
coding across $k$ elements can reduce the per-element bit cost
toward $-\log_2 \fprate$ without integer rounding. The cost is in
query time and implementation complexity: a fractional-bit
fingerprint cannot be compared at machine-word speed and needs the
decoder to reconstruct the integer fingerprint per query. In
practice the $\le 1$ bit per element of waste is acceptable, and
the BHF is deployed with integer $r$. The exercises return to the
question of how much of the wasted bit can be recovered in
worst-case deployments where the target rate is awkward (for
instance $\fprate = 0.06$ gives $-\log_2 \fprate = 4.06$, so $r =
5$ and the achieved FPR is $2^{-5} = 0.03125$, about half the
target). The construction is asymptotically optimal up to the
integer rounding; the rounding itself is small and well-understood,
and the chapter declares the BHF the constructive answer to the
chapter 8 floor question on those terms.
